

\begin{exam} \label{example:thm2blockillustration}
Here we provide examples of the manipulations of the matching field tableaux used throughout the proof of Theorem~\ref{Theorem:2block}. 
The examples we consider are in $\Gr(3,8)$ with the block matching field given by $I_1 = \{1,2,3,4 \}$ and $I_2 = \{5,6,7,8\}$.


The following is an example of a tableau in the semi-standard form as described in the beginning of the proof of Theorem~\ref{Theorem:2block}, 
\[
T = 
\begin{tabular}{ccccc}
 1 & 2 & 5 & 5 & 6 \\
 3 & 3 & 6 & 2 & 4 \\
 4 & 7 & 8 & 7 & 8 \\
 \multicolumn{3}{c}{\upbracefill} & \multicolumn{2}{c}{\upbracefill}\\
 \multicolumn{3}{c}{$C$} & \multicolumn{2}{c}{$D$}\\
\end{tabular}.
\]
The three columns on the left belong to $C$, since the entries are in ascending order and the two columns on the right belong to $D$. The degree of $T$ with respect to the grading described in Lemma~\ref{lem:homo} $(1,1,2,1)$ which can be seen by calculating $\alpha = |\{(1,3,4)\}|$, $\beta = |\{(2,3,7)\}|$ and $\gamma = \{(5,2,7), (6,4,8)\}$. 

%\medskip

As an example of Case~\ref{Case1} in the proof, consider the tableaux
\[
T = 
\begin{tabular}{ccc}
 1 & 3 & 5 \\
 2 & 4 & 3 \\
 4 & 7 & 6 \\
 \multicolumn{1}{c}{\upbracefill} & \multicolumn{2}{c}{\upbracefill}\\
 \multicolumn{1}{c}{$C$} & \multicolumn{2}{c}{$D$}\\
\end{tabular}, 
\quad
T' = 
\begin{tabular}{ccc}
 1 & 3 & 5 \\
 3 & 4 & 2 \\
 4 & 6 & 7 \\
 \multicolumn{1}{c}{\upbracefill} & \multicolumn{2}{c}{\upbracefill}\\
 \multicolumn{1}{c}{$C$} & \multicolumn{2}{c}{$D$}\\
\end{tabular}.
\]
Then we may swap the entries $2$ and $3$ in the second row of $T'$ to obtain a row-wise equal tableau whose left most column is $(1,2,4)$. 

%\medskip
As an example of a tableau manipulation performed in Case~\ref{Case2} of the theorem consider the tableaux, 
\[
T = 
\begin{tabular}{ccc}
 1 & 5 & 6 \\
 2 & 1 & 3 \\
 8 & 6 & 7 \\
 \multicolumn{1}{c}{\upbracefill} & \multicolumn{2}{c}{\upbracefill}\\
 \multicolumn{1}{c}{$C$} & \multicolumn{2}{c}{$D$}\\
\end{tabular},
\quad
T' = 
\begin{tabular}{ccc}
 1 & 5 & 6 \\
 3 & 1 & 2 \\
 6 & 7 & 8 \\
 \multicolumn{1}{c}{\upbracefill} & \multicolumn{2}{c}{\upbracefill}\\
 \multicolumn{1}{c}{$C$} & \multicolumn{2}{c}{$D$}\\
\end{tabular}.
\]
First we can swap entries $2$ and $3$ in the second row of $T'$ to obtain tableau whose left most column is $(1,2,6)$. Then we swap entries $8$ and $6$ in the third row of $T$ to obtain a tableau whose left most column is also $(1,2,6)$. Note that we cannot swap entries $6$ and $8$ in the third row of $T'$.

\medskip
We omit examples of Case~\ref{Case3} in the proof since no swaps occur in this case. For an example of the manipulations made in Case~\ref{Case4}, we consider instead the block matching field 
$I_1 = \{1, 2\}$ and $I_2 = \{3,4,5,6,7,8 \}$. Consider then the two tableaux, 
\[
T = 
\begin{tabular}{cccc}
 4 & 5 & 3 & 6 \\
 5 & 7 & 1 & 2 \\
 6 & 8 & 5 & 7 \\
 \multicolumn{2}{c}{\upbracefill} & \multicolumn{2}{c}{\upbracefill}\\
 \multicolumn{2}{c}{$C$} & \multicolumn{2}{c}{$D$}\\
\end{tabular},
\quad
T' = 
\begin{tabular}{cccc}
 3 & 6 & 4 & 5 \\
 5 & 7 & 1 & 2 \\
 7 & 8 & 5 & 6 \\
 \multicolumn{2}{c}{\upbracefill} & \multicolumn{2}{c}{\upbracefill}\\
 \multicolumn{2}{c}{$C$} & \multicolumn{2}{c}{$D$}\\
\end{tabular}.
\]
First we swap entries $3$ and $4$ in the first row of $T'$ to obtain a tableau whose left most column is $(4,5,7)$. Then we may swap entries $6$ and $7$ in the last row of $T'$ to obtain a tableau whose left most column is $(4,5,6)$ which is identical to the left most column of $T$. Note that we can perform the swap of entries $3$ and $4$ in $T$ however we cannot swap entries $6$ and $7$ in $T$. 

\end{exam}
